

\section{Concerning the Lutheran Free Church and Its Tasks}

\subsection{What the Free Church Needs Most}

Installment 1: “Folkebladet,” February 23, 1916

Most of us are surely aware that, if we are to make some progress toward the accomplishment of our task as a Free Church, the labor required for this will be both long and arduous.

I believe we do well to be clear that the Free Church is still in its formative stage. The goal is still far off. But even though we see the goal only from afar, we must not lose sight of it; rather, we must keep our gaze fixed upon it and continually set our course toward it. The conditions under which we labor to build a Free Church are by no means easy. And by this I am not thinking only of the fact that we are hemmed in between large, strong church bodies, with money and power and with long-established legitimacy and standing among our people.

I am also thinking of the materialistic spirit by which our age is so deeply marked. But if the calling and the task have been entrusted to us by the Lord, we have only one thing to do: to labor tirelessly and steadfastly toward the accomplishment of the task, whether the work seems easy or hard to us—according to the word of Scripture: “Let us do what is good and not grow weary; for we shall reap in due time, provided that we do not lose heart.”

For those who have grown to love the Lutheran Free Church in this country, and who therefore seek to take part in its work, it seems natural to me to ask: What is it that we especially need in order to labor purposefully toward the accomplishment of our task? And how can we, if possible, obtain what we need?

I know that many kinds of answers to these questions present themselves to our minds. It is natural, for example, to think of a seminary, much money, great and gifted leaders, and so forth.

Well, yes, it is true: we can never accomplish anything as a Free Church without a seminary, money, and men to lead the work.

But it is equally true that these and many other necessary things are not enough.

And what else—what more—do we need? There is a word from God that in former times was often heard among us, but which in more recent times we may have been in danger of forgetting: “Not by might, not by power, but by my Spirit, says the Lord.”

I will therefore answer the question of what we need most in this way: We need the Spirit—the Spirit of God. It is true that we also need much else, such as brotherly love, a sense of responsibility, and willingness to work.

But first and foremost, the Spirit of God.

We need the Spirit of God to bear witness to us that Jesus is our life—the life of the individual and of the congregation; and we need his guidance both to see the way and to receive strength to walk in it.

We need the help of the Spirit of God, through the word of Scripture, to come to a clear understanding and full conviction that the position of the Free Church is biblical, and that the truths upon which it is built are the very word of Scripture. And the word of Scripture, despite every attack, stands unshaken.

In order to continue with confidence the church work that has been begun, we need full assurance that, so long as we are willing to walk the straight path—in accordance with the conviction and understanding we have gained from the New Testament’s presentation of congregational life—the Free Church too is unconquerable, despite all public opinion and all manifest enthusiasm for association.

I have not, however, meant by this to assert that we who today constitute the Lutheran Free Church are unconquerable. No, I am fully aware of our weakness and vulnerability, and that it may not be so difficult to cause us to drift in the struggle we have taken up, which many consider both unnecessary and hopeless.

But my thought is this: if we who now constitute the Free Church fail in our calling and our task, given to us by the Lord, and thus forfeit the blessing he has intended for us through this work, he will see to it that others arise, equipped with the light and power needed for the Free Church principle to be realized more and more, as God’s people gather around the free congregation as the right form of God’s kingdom on earth.

By this I do not mean to say that the ordering and work of denominational organizations have no value or legitimacy. On the contrary. I willingly acknowledge that within the church bodies there are congregations with believing members, and I also willingly acknowledge that a church body, by reason of its strong cohesion, has many favorable conditions for carrying out vigorous work, at least outwardly.

I mean only to say this: that we who belong to the association called the Lutheran Free Church and have learned to love it, belong there because we believe that the Lord, through the ecclesiastical development and ecclesiastical conflicts among Norwegian seekers after a church home in this country, has led us there, and because we believe that the New Testament points in that direction: free congregations joined together by “the unity of the Spirit in the bond of peace.”

It is therefore also of exceedingly great importance for us that, under the guidance of the Spirit of God, we receive light concerning the greatness and beauty of our task—the bringing of the congregations to life and their liberation—and receive strength to walk the biblical path in our work of building congregations.

We need, in addition—especially in the Free Church, where no outward bonds exist to hold us together—brotherly love, so that we may be bound together as one, not by factional fences and constitution, not by formulas and laws, but by the strong, holy bond of brotherly love, so that it may be recognized and felt, both among ourselves and by others, that we are a family of brothers and sisters.

And we need a sense of responsibility, so that all of us—the individual friends of the Free Church and the individual congregations—continually live and labor under the consciousness that the Lutheran Free Church’s welfare, progress, or decline depends not first and foremost upon the relation of others to it, but upon my own relation to it—the relation of each individual—and upon the relation of my—of our—congregation. Without some measure of such recognition and awareness, and the sense of responsibility that follows from it, it is surely hopeless to expect any future for the Free Church. And finally—and not least—we need willingness to work, so that we are willing to do something to build our church home.

We must, so far as possible, all take part in this work, so that each and every one of us, according to ability and gifts, endeavors to fill the place in which the Lord has set us, and there to do the work for him with eagerness, willingness, and joy.
